\section*{第 \ref{ch:b1:hotspots} 章——热点与板内火山活动}

\begin{solution}{exo:b1:hotspots:1}
(a) $4205 + 5274 = \qty{9479}{m}$，约 \qty{9.5}{km}。火山
比它所取代的水更重：板块在其载荷下向下弯曲，因此其基底
比远处的海底更深，其真实高度更大。
(b) $74000/0.1 = \num{740000}$ 年（按基拉韦厄的供应速率）；
$74000/0.025 = \num{3.0e6}$ 年（按该海山链的平均速率）。(c)
$0.12\times19 = \qty{2.3}{km^3}$，约为茂纳罗亚火山的 $\num{3e-5}$，即
百分之一的千分之三。
\end{solution}

\begin{solution}{exo:b1:hotspots:2}
(a) $0.92\times\pi/180 = \qty{0.0161}{rad.Ma^{-1}}$。(b) 以弧度每百万年
乘以千米计，$\omega R\sin\Delta$ 的单位为千米每百万年，
即毫米每年：$0.016057\times6371 = 102.3$，因此
$v = 102.3\sin\Delta$，由此得在 \ang{30} 处为 \qty{51.1}{mm.\text{年}^{-1}}，
在 \ang{60} 处为 \qty{88.6}{mm.\text{年}^{-1}}，在 \ang{90} 处为 \qty{102.3}{mm.\text{年}^{-1}}，且
在 \ang{108.0} 处为 \qty{97.3}{mm.\text{年}^{-1}}。(c) 速率与
到轴的距离 $R\sin\Delta$ 成正比，在与极点相距 \ang{90}
的\omterm{def:b1:spherical-astronomy:spherical-triangle}{大圆}上最大；越过此圆后，该点又趋近于轴，
位于对极的一侧。
\end{solution}

\begin{solution}{exo:b1:hotspots:3}
(a) 德干：$\num{e6}/\num{5e5} = \qty{2}{km}$。中大西洋：
$\num{2e6}/\num{7e6} = \qty{0.29}{km}$，约 \qty{290}{m}。(b) 德干的平均值
接近其 \qty{2.5}{km} 的剖面：该省是一片厚实、
连续的高原。(c) 中大西洋
面积即整个区域的面积，该区域横跨四大洲，其残迹
遍布其上：熔岩、岩墙和岩床如今已被大洋分隔，并经受了大量
侵蚀。岩石仅集中于某些地点，在该区域的大部分地方并不存在，
因此平均值无法反映单个堆积体的情况。
\end{solution}

\begin{solution}{exo:b1:hotspots:4}
(a) $\rho_m\alpha\,\Delta T = 3300\times\num{3.2e-5}\times250 =
\qty{26.4}{kg.m^{-3}}$。(b) 隆起柱多出的质量
$(\rho_m - \rho_w)h = 2300\times690 = \qty{1.587e6}{kg.m^{-2}}$，须由
轻质层平衡：$H = \num{1.587e6}/26.4 = \qty{6.0e4}{m}$，约
\qty{60}{km} 的热岩石。
\end{solution}

\begin{solution}{exo:b1:hotspots:5}
(a) 在 A 点（$\lambda = 0$，$\varphi = 0$），$\vec e = (0, 1, 0)$ 且
$\vec n = (0, 0, 1)$；其中 $\sin\ang{270} = -1$，$\cos\ang{270} = 0$，
则 $\vec t_A = -\vec e = (0, -1, 0)$。在 B 点（$\lambda = \ang{90}$），
$\vec e = (-1, 0, 0)$，$\vec n = (0, 0, 1)$，且 $\vec t_B = (1, 0, 0)$。
(b) $\vec t_A\times\vec t_B = (0, 0, 1)$。(c) 取 $\vec p' = (0, 0, 1)$ 且
$\vec r_A = (1, 0, 0)$，$\vec p'\times\vec r_A = (0, 1, 0)$，
$(\vec p'\times\vec r_A)\cdot\vec t_A = -1 < 0$，故该极为
$-\vec p' = (0, 0, -1)$，即南极。沿赤道向西的运动
从北极上方看为顺时针，故从南极上方看为
逆时针。(d) 在距极点 \ang{90} 处，$\omega = 50/6371 =
\qty{7.85e-3}{rad.Ma^{-1}} = \qty{0.45}{\degree\per\Ma}$.
\end{solution}

\begin{solution}{exo:b1:hotspots:6}
(a) $\sum x_i v_i = 90000 + 240000 + 342000 = \num{672000}$，且
$\sum x_i^2 = \num{7.0e7}$，故 $\omega = \qty{0.0096}{rad.Ma^{-1}}$，即
\qty{0.55}{\degree\per\Ma}。(b) 预测值为 28.8、48.0 和
\qty{57.6}{mm.\text{年}^{-1}}；残差分别为 $+1.2$、$0.0$ 和 $-0.6$
\unit{mm.\text{年}^{-1}}。(c) 对于刚性旋转，速率与
到轴的距离成正比，轴上的点不动，因此
不存在需要拟合的截距。
\end{solution}

\begin{solution}{exo:b1:hotspots:7}
(a) $v = \qty{50e-3}{m}/\qty{3.156e7}{s} = \qty{1.58e-9}{m.s^{-1}}$，故
$B = 2300\times\num{3e8}\times\num{1.58e-9} = \qty{1090}{kg.s^{-1}}$，且
$Q = (910/\num{3.2e-5})\times1090 = \qty{3.1e10}{W}$。(b) 约为
夏威夷\omterm{def:b1:heat-budget:heat-flow}{热流}量的三分之一。(c) $B$ 随 $A\,v$ 增大：位于较快板块下方的\omterm{def:b1:hotspots:plume}{地幔柱}
必须以两倍速度重新造出同样的隆起，其\omterm{def:b1:heat-budget:heat-flow}{热流}量也增至两倍。
\end{solution}

\begin{solution}{exo:b1:hotspots:8}
(a) $350/(16.1 - 10.3) = 350/5.8 = \qty{60}{km.Ma^{-1}}$，即
\qty{60}{mm.\text{年}^{-1}}。(b) 年轻的\omterm{def:b1:hotspots:track}{热点轨迹}，\qty{29 \pm 5}{mm.\text{年}^{-1}}，与
板块的 \qty{26 \pm 5}{mm.\text{年}^{-1}} 在不确定度范围内一致；
老段的速率超过两倍，$60 - 26 = 34$，几乎相当于板块速率的七个
标准不确定度。(c) 是的：北偏东 \ang{54} 即\omterm{def:b1:spherical-astronomy:horizontal}{方位角}
\ang{54}，与 \ang{235} 相差不到一度。火山活动在板块上的
运动方向与板块在\omterm{def:b1:hotspots:hotspot}{热点}上方的运动方向相反。(d) 拉伸使较老的破火山口
彼此远离，使它们之间的距离增大，从而
表观速率随之增大；板块本身并未加速。
\end{solution}

\begin{solution}{exo:b1:hotspots:9}
(a) $\num{9.2e10}/\num{e12} = \qty{0.092}{W.m^{-2}}$，即 \qty{92}{mW.m^{-2}}。
(b) 为观测到的过剩\omterm{def:b1:heat-budget:heat-flow}{热流}的九至十八倍，而该过剩\omterm{def:b1:heat-budget:heat-flow}{热流}甚至可能为零。
(c) $L^2/\kappa =
\num{e10}/\num{e-6} = \qty{e16}{s}$，约 \qty{320}{Ma}；板块穿越
\qty{1000}{km} 需 $1000/92 = \qty{11}{Ma}$。来自下方的热量有
时间仅穿过板块的一小部分（此时板块位于隆起之上），因此大部分
\omterm{def:b1:hotspots:plume}{地幔柱}的热量仍留在板块下方的\omterm{def:b1:layered-earth:layers}{地幔}中，并随板块
一同被带走，地表几乎感受不到。
\end{solution}

\begin{solution}{exo:b1:hotspots:10}
(a) 熔体量为十倍，即 \qty{e7}{km^3}，半径为
$(3V/4\pi)^{1/3} = \qty{134}{km}$。(b) $\num{e7}/100 = \qty{e5}{km^2}$，一个圆盘
直径为 $2\sqrt{\num{e5}/\pi} = \qty{357}{km}$，面积比
Traps 小十八倍。\omterm{def:b1:hotspots:plume}{地幔柱}柱头或更大、熔融较少，或铺展得更宽、更薄，
或熔岩流远远超出其范围。(c) 每年 \qty{1}{km^3}，为
基拉韦厄现今供应量的十倍。(d) 新的\omterm{def:b1:hotspots:plume}{地幔柱}必须聚集足够多
炽热而轻质的岩石，以穿过黏性\omterm{def:b1:layered-earth:layers}{地幔}向上推进，而大体积
上升远快于小体积；随后的尾部沿着
柱头开辟的炽热而薄弱的通道向上流动，一条狭窄的通道即可。
\end{solution}

\begin{solution}{exo:b1:hotspots:11}
(a) $\qty{40}{km.Ma^{-1}}\times(81 - 47)\,\unit{Ma} = \qty{1360}{km}$。
(b) 板块以 $(-80\sin\ang{20}, 80\cos\ang{20}) = (-27.4, 75.2)$ 的速度运动，
\omterm{def:b1:hotspots:hotspot}{热点}的速度为 $(0, -40)$：相对于\omterm{def:b1:hotspots:hotspot}{热点}，板块的运动为
$(-27.4, 115.2)$ \unit{mm.\text{年}^{-1}}，速度为 \qty{118.4}{mm.\text{年}^{-1}}，方向
$\ang{360} - \arctan(27.4/115.2) = \ang{346.6}$。(c) 若由轨迹推算，
则速度偏高 \qty{48}{\percent}，方向偏东 \ang{6.6}：
漂移的\omterm{def:b1:hotspots:hotspot}{热点}记录的是板块相对于其自身的运动，而非相对于
深部\omterm{def:b1:layered-earth:layers}{地幔}的运动。
\end{solution}

\begin{solution}{exo:b1:hotspots:12}
(a) $\qty{3e6}{J.m^{-3}.K^{-1}}\times\qty{1300}{K}\times\qty{e15}{m^3} =
\qty{3.9e24}{J}$。(b) 在 $\qty{3.16e13}{s}$ 内，功率为 \qty{1.2e11}{W}，约等于
今日夏威夷\omterm{def:b1:hotspots:plume}{地幔柱}所携热量的 1.3 倍，并占
地球 \qty{47}{TW} 的 \qty{0.26}{\percent}。(c) 这是下限：柱头还带来
其中九成不熔岩石的多余热量，并使
\omterm{def:b1:layered-earth:lithosphere}{岩石圈}底部升温，且潜热未计入。
\end{solution}

\begin{solution}{pb:b1:hotspots:1}
\textbf{1.} 在夏威夷（$\varphi = \ang{19.0}$，$\lambda = \ang{-155.2}$）处，
取 $\sin\ang{304} = -0.8290$ 与 $\cos\ang{304} = 0.5592$：
\[
\begin{aligned}
  \vec r_H &= (-0.8583,\ -0.3966,\ 0.3256), &
  \vec e &= (0.4195,\ -0.9078,\ 0),\\
  \vec n &= (0.2955,\ 0.1366,\ 0.9455), &
  \vec t_H &= (-0.1825,\ 0.8289,\ 0.5287).
\end{aligned}
\]

\textbf{2.} 在路易斯维尔（$\varphi = \ang{-53.6}$，$\lambda = \ang{-140.6}$）处，
取 $\sin\ang{316} = -0.6947$ 与 $\cos\ang{316} = 0.7193$：
\[
\begin{aligned}
  \vec r_L &= (-0.4586,\ -0.3767,\ -0.8049), &
  \vec e &= (0.6347,\ -0.7727,\ 0),\\
  \vec n &= (-0.6220,\ -0.5109,\ 0.5934), &
  \vec t_L &= (-0.8883,\ 0.1693,\ 0.4269).
\end{aligned}
\]

\textbf{3.} $\lVert\vec t_H\rVert^2 = 0.0333 + 0.6871 + 0.2795 = 0.9999$ 及
$\vec t_H\cdot\vec r_H = 0.1566 - 0.3287 + 0.1721 = 0.0000$，在舍入误差范围内。

\textbf{4.} $\vec t_H\times\vec t_L = (0.2643, -0.3917, 0.7054)$，其长度为
$0.8491$。两条轨迹之间的夹角为 $\arcsin0.8491 = \ang{58.1}$（其
点积 $0.528$ 为正，故夹角为锐角）。

\textbf{5.} $\vec p' = (0.3113, -0.4613, 0.8308)$。于是
$(\vec p'\times\vec r_H)\cdot\vec t_H = -0.98 < 0$：极点为
$\vec p = (-0.3113, 0.4613, -0.8308)$。

\textbf{6.} $\arcsin(-0.8308) = \ang{-56.18}$ 及
$\operatorname{atan2}(0.4613, -0.3113) = \ang{124.01}$：极点位于
南纬 \ang{56.2}，东经 \ang{124.0}，位于澳大利亚以南的南大洋。

\textbf{7.} $\vec p\cdot\vec r_H = 0.2672 - 0.1830 - 0.2705 = -0.1863$，故
$\Delta_H = \ang{100.7}$; $\vec p\cdot\vec r_L = 0.1428 - 0.1738 + 0.6687 =
0.6377$，故 $\Delta_L = \ang{50.4}$。

\textbf{8.} 矢量 $\vec p\times\vec r_L$ 垂直于 $\vec p$ 及
$\vec r_L$，$\vec t_L$ 亦如此（由 $\vec p$ 的构造，且因为
轨迹水平）；二者位于同一直线上，其长度为
$\sin\Delta_L$。于是 $(\vec p\times\vec r_L)\cdot\vec t_L = \pm\sin\Delta_L$，且
计算给出 $+0.770 = \sin\ang{50.4}$：板块在路易斯维尔处
沿 $+\vec t_L$ 方向运动，即轨迹延伸的方向。符号法则用于夏威夷时，
本会在路易斯维尔给出错误的方向；但它并未如此，故两者
方向被精确复现，分别为 \ang{304} 与 \ang{316}。

\textbf{9.} $\omega = 67/(6371\sin\ang{50.38}) = 67/4908 =
\qty{0.01365}{rad.Ma^{-1}}$，即 \qty{0.782}{\degree\per\Ma}。

\textbf{10.} $360/0.782 = \qty{460}{Ma}$。

\textbf{11.} 由余弦定理，南纬 \ang{56.18}，东经 \ang{124.01} 与
南纬 \ang{63.78}，东经 \ang{112.48} 两点的夹角为 \ang{9.5}。双轨迹速率，
\qty{0.78}{\degree\per\Ma}，比由十三个速度拟合得到的
\qty{0.92}{\degree\per\Ma} 低 \qty{15}{\percent}。

\textbf{12.} 社会群岛 \ang{73.6}，皮特凯恩 \ang{77.8}，科布 \ang{134.7}。

\textbf{13.} $\omega R = 0.01365\times6371 = \qty{86.98}{mm.\text{年}^{-1}}$，乘以
$\sin\Delta$：社会群岛 \qty{83.5}{mm.\text{年}^{-1}}，皮特凯恩
\qty{85.0}{mm.\text{年}^{-1}}，科布 \qty{61.8}{mm.\text{年}^{-1}}。

\textbf{14.} 皮特凯恩相符：比实测值小 \qty{5}{mm.\text{年}^{-1}}，在其
\qty{15}{mm.\text{年}^{-1}} 的不确定度之内。社会群岛偏慢 \qty{25.5}{mm.\text{年}^{-1}}，为 2.6
个不确定度；科布偏快 \qty{18.8}{mm.\text{年}^{-1}}，超过六个：
两者皆不相符。

\textbf{15.} 极点垂直于两条轨迹矢量。若二者
近乎平行，则其叉积很短，而任一方向上几度的误差
都会使其转过很大的角度：极点
确定得很差。

\textbf{16.} $v = \qty{92e-3}{m}/\qty{3.156e7}{s} = \qty{2.92e-9}{m.s^{-1}}$
则 $B = 2300\times\num{4.83e8}\times\num{2.92e-9} = \qty{3240}{kg.s^{-1}}$。

\textbf{17.} $Q = (910/\num{3.2e-5})\times3240 = \qty{9.2e10}{W}$,
即 \qty{47}{TW} 的 \qty{0.20}{\percent}。

\textbf{18.} 洋中脊每年每千米生成 $21/67000 = \qty{3.1e-4}{km^3}$
；$0.1/\num{3.1e-4} = \qty{320}{km}$ 的洋中脊。

\textbf{19.} $74000/0.1 = \num{740000}$ 年。

\textbf{20.} $\num{2e6}/\num{81e6} = \qty{0.025}{km^3}$ 每年：仅基拉韦厄
现今所获的供应量为整个\omterm{def:b1:hotspots:hotspot}{热点}平均产量的四倍，而该\omterm{def:b1:hotspots:hotspot}{热点}
正处于接近其最强的时期。

\textbf{21.} $v_H = 86.98\sin\ang{100.73} = 86.98\times0.9825 =
\qty{85.5}{mm.\text{年}^{-1}}$。

\textbf{22.} $92 - 85.5 = \qty{6.5}{mm.\text{年}^{-1}}$，为实测速度的 2.2 个标准
不确定度：两个\omterm{def:b1:hotspots:hotspot}{热点}对于固定参考系中一个
\omterm{def:b1:plate-kinematics:rigid-plate}{刚性板块}的运动，并不完全相符。

\textbf{23.} 由年龄得出，每年 \qty{9.1}{cm} 与 \qty{8.6}{cm}，均
对整条海山链取平均；摩尔根所得的每年 \qty{9.2}{cm} 接近
前者，路易斯维尔的预测值每年 \qty{8.5}{cm} 接近
后者。四者彼此相差均在每年 \qty{0.7}{cm} 之内：轨迹、年龄
与方向给出的结论，在约 \qty{8}{\percent} 之内一致。

\textbf{24.} 保留一位小数，路易斯维尔轨迹给出太平洋板块相对
于夏威夷\omterm{def:b1:hotspots:hotspot}{热点}的速度为
\textbf{\qty[text-series-to-math=true]{8.5}{cm.\text{年}^{-1}}}。
\end{solution}
